\subsection{结果分析}

\textbf{三分图先验带来的差距。} 由表~\ref{tab:pm} 可见，享有三分图先验的 ViTMatte 在全图口径上明显领先（SAD 25.5、MSE 0.012），因为前景范围已由三分图给定，模型只需在狭窄的未知带上精修；三种无三分图方法则需自主完成前景判定，其全图 SAD 普遍更高（SDMatte 44.3、SMat-auto 64.1、ZIM-L 93.2）。这一差距并非单纯的「模型优劣」，而是输入条件不同的直接体现，正对应「给定前景的精修」与「自主前景选择」两条路线的本质区别。

\textbf{全图与未知区域的分歧。} 值得注意的是，若只看\emph{未知区域}指标，SDMatte（SAD 23.9）甚至略优于 ViTMatte（24.3），说明在边缘细节的局部刻画上，扩散先验与三分图精修旗鼓相当；三种无三分图方法的全图误差主要来自\emph{未知区域之外}的前景范围误判，而非边缘本身。这印证了同时报告两套口径的必要性——全图指标度量「主体识别/图形-背景分离」，未知区域指标度量「边缘细节质量」，二者分离才能定位误差来源。

\textbf{无三分图方法之间的比较。} 三者中扩散式 SDMatte 最稳健，在全部 150 张上零图形-背景反转；SMat 自动模式凭借 DINOv2 显著性前端，零提示即可避免反转，交互（8 点）相比自动几乎无提升，说明其自动分支已抓对前景。相比之下，ZIM 的软 matte 头在「主体满框」的植株上常把更大、更干净的背景区误判为前景，即使给定真值框与 8 个前景点仍会发生图形-背景反转，是其全图误差显著偏高（SAD 93.2）的主因。这说明在植物域，自动前端的前景判定（而非边缘刻画）才是主要瓶颈。

\textbf{定性观察。} 图~\ref{fig:qual} 直观印证上述结论：在典型样本（如第 1 行）上各模型预测均贴近真值；而在主体满框、与背景同种植株相似的难例（第 3、4 行）上，无三分图方法出现前景范围扩张或细枝丢失，ViTMatte 则因三分图约束始终稳定。结合跨数据集自验（表~\ref{tab:crossval}）中各模型在通用基准上的良好表现可知，植物场景的困难主要来自数据域本身的图形-背景模糊性，而非模型在边缘刻画上的不足。
